% ============================================================
% Prime — الصفحة الأولى: صندوق إعادة الإنتاج
% نسخة الإيداع النهائية — تُستبدل الروابط النائبة بالروابط الفعلية فقط
% لا تستخدم هذه الصيغة قبل إغلاق GitHub Release + Zenodo
% ============================================================

\begin{tcolorbox}[
  colback=gray!5,
  colframe=black!70,
  title={\textbf{إعادة الإنتاج والمواد الحسابية (Reproducibility and Computational Materials)}},
  fonttitle=\bfseries,
  breakable,
  left=4mm, right=4mm, top=3mm, bottom=3mm
]

تُتاح الأكواد الحسابية ومواد معالجة البيانات المستخدمة في هذه الدراسة للتحقق من النتائج وإعادة إنتاج الحسابات.

\vspace{1mm}
\noindent\textbf{مستودع الأكواد (GitHub):} \url{[Repository Link]}

\vspace{1mm}
\noindent\textbf{الإصدار البرمجي المعتمد:} \url{[Release Link]}

\vspace{1mm}
\noindent\textbf{أرشيف الإصدار (Zenodo):} \url{[Code DOI]}

\vspace{1mm}
يرتبط هذا الإصدار بالحسابات والنتائج الواردة في الدراسة، ويُتيح الاطلاع على خطوات المعالجة وإعادة تنفيذها وفق الإعدادات الموثقة.

\end{tcolorbox}

% ------------------------------------------------------------
% تعليمات الاستبدال عند الإيداع — لا تظهر في الورقة:
% [Repository Link] = رابط مستودع GitHub
% [Release Link]   = رابط GitHub Release المحدد، مثال: v1.0-prime
% [Code DOI]       = رابط DOI الخاص بأرشيف Zenodo لنفس الإصدار
% شرط الإغلاق: Release و Zenodo يجب أن يمثلا النسخة نفسها تمامًا.
% ------------------------------------------------------------
